\section{Local rings}
\label{section-local-rings}

\noindent
Local rings are the bread and butter of algebraic geometry.

\begin{definition}
\label{definition-local-ring}
A {\it local ring} is a ring with exactly one maximal ideal.
If $R$ is a local ring, then the maximal ideal is often denoted
$\mathfrak m_R$ and the field $R/\mathfrak m_R$ is called the
{\it residue field} of the local ring $R$.
We often say ``let $(R, \mathfrak m)$ be a local ring''
or ``let $(R, \mathfrak m, \kappa)$ be a local ring''
to indicate that $R$ is local, $\mathfrak m$ is its unique
maximal ideal and $\kappa = R/\mathfrak m$ is its residue field.
A {\it local homomorphism of local rings} is a ring map
$\varphi : R \to S$ such that $R$ and $S$ are local rings and such
that $\varphi(\mathfrak m_R) \subset \mathfrak m_S$.
If it is given that $R$ and $S$ are local rings, then the phrase
``{\it local ring map $\varphi : R \to S$}'' means that $\varphi$
is a local homomorphism of local rings.
\end{definition}

\noindent
A field is a local ring. Any ring map between fields is a local
homomorphism of local rings.

\medskip\noindent
The localization $R_\mathfrak p$ of a ring $R$ at a prime $\mathfrak p$
is a local ring with maximal ideal $\mathfrak p R_\mathfrak p$. Namely, by
Lemma \ref{lemma-spec-localization} every prime ideal of $R_\mathfrak p$
is contained in the prime ideal $\mathfrak p R_\mathfrak p$
(hence this is a maximal ideal and the only maximal ideal of $R_\mathfrak p$).
The residue field of $R_\mathfrak p$ is denoted $\kappa(\mathfrak p)$;
we call it the {\it residue field of $\mathfrak p$}; by
Proposition \ref{proposition-localize-quotient} we may identify
$\kappa(\mathfrak p)$ with the field of fractions of the domain $R/\mathfrak p$.
Via the composition
$$
\Spec(\kappa(\mathfrak p)) \to \Spec(R_\mathfrak p) \to \Spec(R)
$$
the unique point of the source maps to the point $\mathfrak p$ of the target.

\medskip\noindent
Let $\varphi : R \to S$ be a ring map. Let $\mathfrak q \subset S$ be a prime
and consider the prime $\mathfrak p = \varphi^{-1}(\mathfrak q)$ of $R$.
Since $\varphi(\mathfrak p) \subset \mathfrak q$ the induced ring map
$$
R_\mathfrak p \to S_\mathfrak q,\quad r/g \mapsto \varphi(r)/\varphi(g)
$$
is a local ring map and we obtain an induced map of residue fields
$\kappa(\mathfrak p) \to \kappa(\mathfrak q)$.


\begin{example}
\label{example-not-local}
If $R$ is a local ring and $\mathfrak p \subset R$ is a non-maximal
prime ideal, then $R \to R_\mathfrak p$ is not a local homomorphism.
\end{example}

\begin{lemma}
\label{lemma-characterize-local-ring}
Let $R$ be a ring. The following are equivalent:
\begin{enumerate}
\item $R$ is a local ring,
\item $\Spec(R)$ has exactly one closed point,
\item $R$ has a maximal ideal $\mathfrak m$
and every element of $R \setminus \mathfrak m$
is a unit, and
\item $R$ is not the zero ring and for every $x \in R$ either $x$
or $1 - x$ is invertible or both.
\end{enumerate}
\end{lemma}

\begin{proof}
Let $R$ be a ring, and $\mathfrak m$ a maximal ideal.
If $x \in R \setminus \mathfrak m$, and $x$ is not a unit
then there is a maximal ideal $\mathfrak m'$ containing $x$.
Hence $R$ has at least two maximal ideals. Conversely,
if $\mathfrak m'$ is another maximal ideal, then choose
$x \in \mathfrak m'$, $x \not \in \mathfrak m$. Clearly
$x$ is not a unit. This proves the equivalence of (1) and (3).
The equivalence (1) and (2) is tautological.
If $R$ is local then (4) holds since $x$ is either in $\mathfrak m$
or not. If (4) holds, and $\mathfrak m$, $\mathfrak m'$ are distinct
maximal ideals then we may choose $x \in R$ such that
$x \bmod \mathfrak m' = 0$ and $x \bmod \mathfrak m = 1$
by the Chinese remainder theorem
(Lemma \ref{lemma-chinese-remainder}).
This element $x$ is not invertible and neither is $1 - x$ which is
a contradiction. Thus (4) and (1) are equivalent.
\end{proof}

\begin{lemma}
\label{lemma-characterize-local-ring-map}
Let $\varphi : R \to S$ be a ring map. Assume $R$ and $S$ are local rings.
The following are equivalent:
\begin{enumerate}
\item $\varphi$ is a local ring map,
\item $\varphi(\mathfrak m_R) \subset \mathfrak m_S$,
\item $\varphi^{-1}(\mathfrak m_S) = \mathfrak m_R$, and
\item for any $x \in R$, if $\varphi(x)$ is invertible in $S$, then $x$
is invertible in $R$.
\end{enumerate}
\end{lemma}

\begin{proof}
Conditions (1) and (2) are equivalent by definition.
If (3) holds then (2) holds. Conversely, if (2) holds, then
$\varphi^{-1}(\mathfrak m_S)$ is a prime ideal containing
the maximal ideal $\mathfrak m_R$, hence
$\varphi^{-1}(\mathfrak m_S) = \mathfrak m_R$. Finally, (4) is the
contrapositive of (2) by Lemma \ref{lemma-characterize-local-ring}.
\end{proof}

\begin{remark}
\label{remark-fundamental-diagram}
A fundamental commutative diagram associated to a ring map
$\varphi : R \to S$ and a prime $\mathfrak p \subset R$ is the following
$$
\xymatrix{
\kappa(\mathfrak p) \otimes_R S =
S_{\mathfrak p}/{\mathfrak p}S_{\mathfrak p}
&
S_{\mathfrak p} \ar[l]
&
S \ar[r] \ar[l]
&
S/\mathfrak pS \ar[r]
&
(R \setminus \mathfrak p)^{-1}S/\mathfrak pS
\\
\kappa(\mathfrak p) =
R_{\mathfrak p}/{\mathfrak p}R_{\mathfrak p} \ar[u]
&
R_{\mathfrak p} \ar[u] \ar[l]
&
R \ar[u] \ar[r] \ar[l]
&
R/\mathfrak p \ar[u] \ar[r]
&
\kappa(\mathfrak p) \ar[u]
}
$$
In this diagram the outer left and outer right columns are identical.
On spectra the horizontal maps induce homeomorphisms
onto their images and the squares induce fibre squares of topological spaces
(see Lemmas \ref{lemma-spec-localization} and \ref{lemma-spec-closed}).
This shows that $\mathfrak p$ is in the image of the map on Spec if and
only if $S \otimes_R \kappa(\mathfrak p)$ is not the zero ring.
If there does exist a prime $\mathfrak q \subset S$ lying over
$\mathfrak p$, i.e., with $\mathfrak p = \varphi^{-1}(\mathfrak q)$
then we can extend the diagram to the following diagram
$$
\xymatrix{
\kappa(\mathfrak q) = S_{\mathfrak q}/{\mathfrak q}S_{\mathfrak q}
&
S_{\mathfrak q} \ar[l]
&
S \ar[r] \ar[l]
&
S/\mathfrak q \ar[r]
&
\kappa(\mathfrak q)
\\
\kappa(\mathfrak p) \otimes_R S =
S_{\mathfrak p}/{\mathfrak p}S_{\mathfrak p} \ar[u]
&
S_{\mathfrak p} \ar[u] \ar[l]
&
S \ar[u] \ar[r] \ar[l]
&
S/\mathfrak pS \ar[u] \ar[r]
&
(R \setminus \mathfrak p)^{-1}S/\mathfrak pS \ar[u]
\\
\kappa(\mathfrak p) =
R_{\mathfrak p}/{\mathfrak p}R_{\mathfrak p} \ar[u]
&
R_{\mathfrak p} \ar[u] \ar[l]
&
R \ar[u] \ar[r] \ar[l]
&
R/\mathfrak p \ar[u] \ar[r]
&
\kappa(\mathfrak p) \ar[u]
}
$$
In this diagram it is still the case that the outer left and outer right
columns are identical and that on spectra the horizontal maps induce
homeomorphisms onto their image.
\end{remark}

\begin{lemma}
\label{lemma-in-image}
Let $\varphi : R \to S$ be a ring map. Let $\mathfrak p$
be a prime of $R$. The following are equivalent
\begin{enumerate}
\item $\mathfrak p$ is in the image of
$\Spec(S) \to \Spec(R)$,
\item $S \otimes_R \kappa(\mathfrak p) \not = 0$,
\item $S_{\mathfrak p}/\mathfrak p S_{\mathfrak p} \not = 0$,
\item $(S/\mathfrak pS)_{\mathfrak p} \not = 0$, and
\item $\mathfrak p = \varphi^{-1}(\mathfrak pS)$.
\end{enumerate}
\end{lemma}

\begin{proof}
We have already seen the equivalence of the first two
in Remark \ref{remark-fundamental-diagram}. The others
are just reformulations of this.
\end{proof}

\begin{remark}
\label{remark-local-ring-fibre}
Let $R \to S$ be a ring map. Let $\mathfrak q$ be a prime
ideal of $S$ lying over the prime ideal $\mathfrak p$ of $R$.
According to Remark \ref{remark-fundamental-diagram}
the prime $\mathfrak q$ corresponds to a unique prime
$\overline{\mathfrak q}$ of the fibre ring
$F = S \otimes_R \kappa(\mathfrak p)$.
Then we have
$$
F_{\overline{\mathfrak q}}
\cong S_\mathfrak q \otimes_{R_\mathfrak p} \kappa(\mathfrak p)
\cong S_\mathfrak q/\mathfrak p S_\mathfrak q
$$
Namely, there is an obvious ring map
$F \to S_\mathfrak q \otimes_{R_\mathfrak p} \kappa(\mathfrak p)$
which is easily seen to be isomorphic to $F \to F_{\overline{\mathfrak q}}$.
The second isomorphism follows from the fact that $\kappa(\mathfrak p)$
is the quotient of $R_\mathfrak p$ by $\mathfrak pR_\mathfrak p$.
\end{remark}







